\section{Kernelized Lifted-Band Branch}
\label{sec:kernelized-lifted-band-branch}

The kernelized lifted-band branch refines the scale-barrier problem inside the
region between the lower-prefix family and the active shell. Its positive-forward
aim is to prove a separated-band inequality with explicit decay in \(j-k\).
The proof burden is the displayed kernel estimate below; scale separation by
itself is not a payment. The branch asks whether that proved separation gain is
strong enough to pay the remaining lower-shell bill.

The target separated-band estimate has the form
\[
  \big|\langle P_j((P_k u\cdot\nabla)u^{\mathrm{near}}_j),u_j\rangle\big|
  \le C_{\alpha}2^{-\alpha(j-k)} A_{j,k}(t),
\]
where \(k<j-C\), \(\alpha>0\), \(C_\alpha\) is independent of \(j\) and \(k\),
and \(A_{j,k}\) contains the active and lower-shell norms. The desired outcome
is summability over \(k<j-C\). A true lifted-band gain would turn separation
into payment.

The branch is valuable because it isolates one precise source of hope. The
unpaid lower-prefix object is not treated as a black box. It is placed into a
kernel, and the kernel is tested for decay. The branch asks whether the geometry
of frequency separation itself can supply the missing smallness.

The stopping point is the half-derivative remainder. The kernel gain can reduce
the lower-shell interaction, yet after the derivative in the nonlinearity is
accounted for, one is left with an extra active or lower-shell factor that is not
paid by the basic energy surface. In the dyadic bookkeeping for the displayed
kernel gain,
\[
  2^{-(j-k)}\cdot 2^j = 2^k,
\]
so the derivative can reappear at the lower shell. The separation has helped, yet
it has not eliminated the derivative bill. The branch transforms the obstruction
rather than deleting it.

This produces a lifted-band residue. The residue is more specific than a generic
source reserve: it is a lower-shell derivative demand exposed inside a separated
interaction. That precision matters because it tells the later CM-exit analysis
what kind of terminal witness the branch can produce.

A positive-forward proof can try to continue from here in several ways. It can
seek stronger kernel decay, impose active-window smallness, use cancellation,
exploit divergence-free structure, or move the unpaid term into a Volterra memory
term. Each move sharpens the residue. Each move also has to explain what
happens at a terminal time if the lower-shell derivative demand survives.

In the CM-exit program, the lifted-band residue becomes proof material after it
is attached to a selected terminal object. If it enters as a lower-shell packet failure, it
presses the Pack face. If it survives through the same nonlinear participation
law, it presses the Part face. If it appears only as loss of coherent positive
scale readout, it presses the Field face. The branch is a technical source of
terminal witnesses, not a closed positive smoothness proof.

The conclusion is narrow and useful. Kernelization improves the picture by
localizing the unpaid bill. It does not by itself close the cascade. The branch
must be carried forward as a source of classified terminal residues inside the
later class-membership exit proof.
